\begin{theorem} \textbf{Teorema de Green.} \label{green}  Sean $U\subseteq \Rn{2}$ abierto  y   $D \subsetneq  U$ un conjunto  simplemente conexo (\autoref{def:conjuntosimpleconexo}) cuyo borde, $\partial D$, es una curva simple,  cerrada y suave a trozos.  Si  $\mathbf{F}:U\to\Rn{2}$ es de clase $\mathcal{C}^1$, entonces  
    \[
        \oint_{\partial D^{+}} \mathbf{F}\cdot d\mathbf{s} = \iint_D \grad\times\mathbf{F}  \: dA.
    \]
\end{theorem}

\begin{obs} Sea $\mathbf{F}:U\subset \Rn{2}\to\Rn{2}$ bajo las condiciones de (\ref{green})  con $\mathbf{F}=(P,Q)$,  el teorema de Green  puede ser escrito de siguiente manera alternativa:
\[
    \oint_{\partial D^+}P\:dx+Q\:dy=\iint_D\left(\dif{Q}{x}-\dif{P}{y}\right)dA,
\]
basta con ver la notaci\'on dada en (\ref{not_int}) y recordar que 
\[
    \nabla\times\mathbf{F} = \dif{Q}{x}-\dif{P}{y}.
\]
\end{obs}

\begin{example} Evaluar la integral de línea $\oint_C -y \, dx + x \, dy$, donde $C$ es la curva simple y cerrada dada por la ecuaci\'on $x^2 + y^2 = 1$ orientada en sentido antihorario.
\end{example}
\noindent\textbf{Solución.}  Antes de aplicar el Teorema de Green, verificamos que se cumplan las hipótesis fundamentales:

\begin{itemize}
    \item El campo vectorial $\mathbf{F}(x,y) = (-y, x)$ tiene componentes $P(x,y) = -y$ y $Q(x,y) = x$. Al ser polinomios, son funciones de clase $\mathcal{C}^1 (\mathbb{R}^2)$.
    \item La curva $C$  es una curva cerrada, simple, suave y está orientada en sentido antihorario.
    \item La región $D$ encerrada por la curva (el disco de radio 1) es simplemente conexa.
\end{itemize}

Luego, el  Teorema de Green nos permite transformar la integral de línea a lo largo de la curva $C$ en una integral de área  sobre la región $D$ usando la siguiente relación:
\[ \oint_C P \, dx + Q \, dy = \iint_D \left( \frac{\partial Q}{\partial x} - \frac{\partial P}{\partial y} \right) \, dA \]

Calculamos las derivadas parciales de nuestro campo:
\[ \frac{\partial Q}{\partial x}(x,y) = 1 \qquad \text{y} \qquad \frac{\partial P}{\partial y}(x,y) = -1 \]

Sustituimos estos valores para pasar explícitamente a la integral de área:
\[ \oint_C -y \, dx + x \, dy = \iint_D (1 - (-1)) \, dA = \iint_D 2 \, dA =  2  \iint_D  \, dA \]

Usando que $D$ es un disco de radio $R = 1$  (por lo que su área es $\pi \cdot 1^2 = \pi$)  y recordando que,  
\[  \iint_D  \, dA =  \text{Área}(D) \]

Obtenemos finalmente que  
\[ \oint_C -y \, dx + x \, dy = 2\pi \]  

\begin{figure}[h!]
    \centering
    \tikzsetnextfilename{ej_green} 
    \begin{tikzpicture}
        % Ejes
        \draw[->] (-2,0) -- (2,0) node[right] {$x$};
        \draw[->] (0,-2) -- (0,2) node[above] {$y$};
        
        % Región D sombreada (acá el circle común no molesta porque no lleva flecha)
        \filldraw[blue!20, opacity=0.5] (0,0) circle (1.5cm);
        
        % Curva C dibujada como ARCO para que la decoración no tire error
        % Arranca en (1.5,0) y hace un giro completo de 0 a 360 grados con radio 1.5cm
        \draw[thick, blue, postaction={decorate, decoration={
            markings, 
            mark=at position 0.25 with {\arrow{Stealth}}
        }}] (1.5,0) arc (0:360:1.5cm);
        
        % Etiquetas
        \node at (0.5, 0.5) {$D$};
        \node[blue] at (1.2, 1.2) {$C$};
    \end{tikzpicture}
    \caption{Región de integración $D$ y su frontera orientada $C$.}
\end{figure}


\begin{tarea}
 Para  la curva simple, cerrada y orientada $C$ formada por los segmentos  de $(0, 0)$ a $(1, 1),$ de (1,1) a $(0, 1)$ y de $(0, 1)$ vuelta al $(0, 0)$, calcular  \[ \int_{C} xy \; dx + \sqrt{y^{2}+1} \; dy.\] 
\end{tarea}


\begin{tarea}
 Sea    $D=\{  (x,y) \in \mathbb{R}^{2} \: : \: 1 \leq x ^{2}+ y^{2}  \leq  4,  x\geq 0 \}$ , calcular    \[ \int_{(\partial D)^{-}} x^{2}y \; dx - xy^{2} \; dy.\] 
\end{tarea}
